\section{Nietzsche and Being}

\emph{The problem of the forgetting of Being cannot be resolved philosophically or scientifically, but only spiritually and metaphysically.}

We started Gornahoor first of all as an exploration in the recovery of Tradition in the West. Since Tradition is the opposite of modernity, this led to an attempt to forge a new relationship to the modern world:

\begin{itemize}


\item First, to recognize the Crisis of the modern world

\item Second, to revolt against it.
\end{itemize}


Following \textbf{Rene Guenon} (which is also the advice of the Dalai Lama, by the way), the man of Tradition should follow the spiritual path of his fathers. This led to two choices:

\begin{enumerate}


\item The pre-schism Catholic/Orthodox tradition, which formed the Traditional civilization of the Middle Ages.

\item Pre-Christian paganism, which formed the Traditional civilizations of the Ancient Greeks, Romans, and Germanics
\end{enumerate}


Choice (1) was preferred by Guenon. Choice (2), by \textbf{Julius Evola}, although his view of paganism was enhanced by his particular explications of Buddhism, Tantrism, and Hermetism. In the next section, we will show why (1) is the only viable option because:

\begin{itemize}


\item Pre-Christian traditions are unrecoverable at this point because they lack all pretense of continuity

\item Everything valid in (2) is included in (1)
\end{itemize}


Nevertheless, a third ``option'' has also appeared: it includes a mix of neo-paganism and Nietzscheism. It is fundamentally anti-traditional and vague and doomed to deliquescence\footnote{\url{http://www.gornahoor.net/?p=3}} ... Evola himself opposed neo-paganism.


\paragraph{The Historical Situation}


Guenon made the judgment that tradition was totally lost in the West, so he moved away. Although he referred to the Catholic tradition, a fortiori, that applies to any pre-Christian pagan tradition. However, just as a perfect Void cannot manifest, neither can Tradition be fully lost. Through the writings of Guenon and Evola, we can see the lineaments of what an actual Tradition looks like. That outline provides the hermeneutic key to understand the esoteric writing of the West, which we can now interpret and understand. Furthermore, certain spiritual currents have come to light subsequent to Guenon's time, which definitively reveal the continuity of a western tradition.

We also know that both the Ancient and the Medieval western civilizations are representative of Tradition, so any reconstruction of Tradition needs to inspect them carefully. Now few things are more jarring, at least to me, like the transition in the Louvre when one leaves the ancient art area and walks up to the medieval section. Although the artistic sensibilities had certainly changed, there must be a relationship of continuity.

That continuity was known to the medievals as we demonstrated through Dante, Boethius, and Lewis' Discarded Image\footnote{\url{http://www.gornahoor.net/?p=3545}}. We have shown that the medievals regarded the greatest of the pagans as worthy of emulation\footnote{\url{http://www.gornahoor.net/?p=331}}. We have shown, through Augustine, that the Christian religion was just the most recent form of a more ancient Tradition.

Furthermore, the early Fathers regarded the Christian religion as the esoteric teachings of Greek philosophy\footnote{\url{http://www.gornahoor.net/?p=7152}}, considered as exoteric. We have shown that Greek philosophy is more than Plato and Aristotle, who are just links in a longer chain\footnote{\url{http://www.gornahoor.net/?p=1256}}. Finally, the spiritual traditions of Buddhism, Hinduism, Hermetism, etc., that Evola points to, fit comfortably into the Western tradition. Buddha, as St. Josaphat\footnote{\url{http://www.firstthings.com/web-exclusives/2009/09/saint-sakyamuni}}, can be comfortably incorporated. As recent works by Valentin Tomberg and Boris Mouravieff demonstrate, Hermetism is alive and well in our time.

The point is that if anyone wants to reconstruct a western tradition on another basis will nevertheless have to go through the same steps. Most importantly, continuity needs to be established, so the spirituality of the west when it was Traditional needs to be part of that program. That we don't see happening.

\paragraph{The Doctrine of the Two Worlds}


Julius Evola, in the opening paragraphs of Revolt Against the Modern World, established the fundamental difference between the Traditional and the modern mentality: viz., the knowledge of Being. I have taken the liberty to retranslate those paragraphs; it may be less elegant, but it is more metaphysically precise. As you read it, keep this point in mind:

\textbf{To the extent that you don't understand the world of Being, then you are still a modern man.}

\begin{quotationx}
In order to understand both the traditional spirit and modern civilization as its negation, it is necessary to start from a fundamental point: the doctrine of the two worlds.

There is a physical order and a metaphysical order. There is a mortal world and the world of the immortals. There is the higher reason of ``being'' and the lower one of ``becoming''. More broadly: there is the visible and tangible and beyond it, first of all, there is the invisible and intangible as the superworld, origin, and true life.

Everywhere in the world of Tradition, both in the East and the West, in one form or another, this knowledge was always present as an unshakeable axis around which everything else was ordered.

We say knowledge and not ``theory''. As difficult as it is for moderns to conceive it, it is necessary to start from the idea that traditional man knew of the reality of an order of being much vaster than what today corresponds to the word ``real''. Today, fundamentally, he no longer conceives as reality anything that lies beyond the world of bodies in space and time. Certainly, there are those who still admit something beyond the sensible: however, it is always in the form of an hypothesis or a scientific law, a speculative idea, or a religious dogma that he admits this something, so actually he does not surpass that limitation: practically, i.e., as direct experience, whatever difference there is in his ``materialistic'' and ``spiritualistic'' beliefs, the average modern man forms his image of reality only in relation to the world of bodies.
\end{quotationx}


Now the world of Being is what we call the God, as much as Evola prefers to avoid that term. Now Being is not known discursively through theories or dogmas, but intuitively through a direct knowledge or gnosis. Hence any recovery of Tradition must move beyond verbiage and must include a path to this gnosis.

\paragraph{Nietzsche's Dilemma}


Now the counter-traditional attempts to recover tradition in terms of neo-paganism and Friedrich Nietzsche, are faced with an immediate difficulty: \emph{Nietzsche denied the world of Being and claimed that only the world of becoming is real.} In the hands of a Julius Evola, who understood Tradition, any value in Nietzsche can extracted with care. However, lesser minds don't even recognize the problem, never mind try to deal with it.

Nietzsche did not offer arguments against God, but rather announced the fact the God is dead. That is just a description of the psychic state of European man of the 19\textsuperscript{th} century: \emph{He had forgotten Being}. Despite that, the consequences had not yet dawned on him. Specifically, without Being, there is no ground for the world, neither intellectually nor morally, or as Evola put it, ``the origin and true life'' of the world of becoming had been completely forgotten. Moreover, as the corollary, there is no purpose to becoming and no direction to life.

The mass of men continue to live as though that psychic change had never happened, living off the crumbs, as it were, of their more Traditional ancestors. If there is no intellectual or moral foundation to the world of becoming, then all that is left is the Will. All worldviews are merely perspectives and the myriad attempts to impose them on others are simply manifestations of the Will to Power. Only habit allows such men to regard those worldviews are true or false, good or evil.

\textbf{To the extent that you believe that what is necessary is a better scientific theory, a better idea, or a better dogma, you don't understand what Nietzsche was driving at.}

Of course, Evola did understand. Hence, his solution -- and ours -- is to gain a direct, intuitive, unmediated knowledge of Being.


\flrightit{Posted on 2014-10-15 by Cologero}

\begin{center}* * *\end{center}

\begin{footnotesize}
\begin{sffamily}
\texttt{X on 2014-10-16 at 11:44 said:}

Being Nietzsche is against Nietzsche's Being.

\hfill

\texttt{Kenneth Anderson on 2014-10-16 at 16:18 said:}

How Nietzsche could have retained God Philosophy since Nietzsche has often centered on the subject and not the object, finding the object only the subjective fixing of one subject by another subject. That was convenient for describing a relativity of values and morals.

Nietzsche, rightly, said that Being will have to be conceived as the sensation which is no longer based on anything quite devoid of sensation. But then since God had been defined as devoid of sensation, Nietzsche rejected God. If Nietzsche had understood Godhood as a material/supermaterial object, or objects, that we evolve to become in the material world, which had been previously, traditionally, but partially, seen as the Inward God or the Father Within, then Nietzsche might have been able to retain God transformed in the Outward Godhood we can evolve to become, if we are successful in survival and evolution---first glimpsed as the Inward God.

Philosophical naturalism for me, which is inherent in theological materialism, doesn't think only in terms of subjective relativity, human nature is determined in many ways, and other parts of nature are as well, and their definition doesn't strictly depend on who is doing the defining. We see as much of reality as our level of consciousness can see.

Getting back to the object means moving away from subjective definition of objects. Human nature and Godhood are seen as objects in reality, with many, although not all determined features, which are simply not defined as only relative to the subject viewing them.

This describes a living evolving object Godhood. In getting back to the real object and moving away from the relative subject we find real Godhood. We don't require the aristocratic radicalism of Nietzsche's solution---which involved a goalless and relative will-to-power---an evolving conservatism will do.

\hfill

\texttt{Cassiodorus on 2014-10-16 at 18:31 said:}

``Now Being is not known discursively through theories... ''

Does the Aristotelian-Thomistic tradition represent an anomaly to this position?

\hfill

\texttt{Cologero on 2014-10-17 at 09:01 said:}

The Thomist position is that discursive knowledge of God is known only ``analogically''.

\hfill

\texttt{exit on 2014-10-18 at 11:31 said:}

Perhaps you are unaware of the Dalai Lamas new book which calls for a global shift towards a universal secular morality and towards progress.

\hfill

\texttt{Cologero on 2014-10-19 at 06:19 said:}

I don't know which book you are referring to, Exit, but it is not surprising. It is the result of confounding religion with moralism: if you arrive at the moral consensus of the modern mind, then it does not matter how you got there.

\hfill

\texttt{Max on 2014-10-26 at 05:32 said:}

A negation is actually a type of affirmation, although in an unconscious way. Exclaiming that god is dead masks a wish to return. Without first acknowledging the distance is it possible? Without knowing our sins could we be free of them? It is not possible to conceive of an ``isolated'' negation without its positive counterpart. All negations thus carries a hidden truth within, otherwise they would not exist. This has radical consequences because it means that nothing is ever lost and that there is always hope. What happens then when we negate a negation? Say only yes or no without any hesitation. Yes to that which is, and no to nothing. To the degree a person says yes and dispels confusion he has life and consciousness.

\hfill

\end{sffamily}
\end{footnotesize}

